\documentclass[10pt,a4paper]{amsart}
\usepackage[margin=19mm]{geometry}
\usepackage{amsmath,amssymb,mathtools}
\usepackage[scr=rsfs,cal=euler]{mathalfa}
\usepackage{xfrac}
\usepackage[unicode,hidelinks,bookmarks=false]{hyperref}
\newcommand{\ie}{i.e.\ }
\newcommand{\cf}{cf.\ }
\newcommand{\resp}{respectively\ }
\newcommand{\Subsection}{\subsection}
\providecommand{\nobreakdash}{\nobreak\hbox{-}}
\setlength{\emergencystretch}{2em}
\multlinegap0pt
\usepackage{bm}
\usepackage{bbm}
\usepackage{dsfont}
\usepackage{xspace}
\usepackage{cancel}
\usepackage{mathtools}
\usepackage{amsthm}
\makeatletter
\@ifundefined{theorem}{\newtheorem{theorem}{Theorem}}{}
\@ifundefined{proposition}{\newtheorem{proposition}[theorem]{Proposition}}{}
\makeatother
\DeclareMathOperator{\essinf}{ess\, inf}
\DeclareMathOperator{\esssup}{ess\,sup}
\NewDocumentCommand{\dmu}{}{\ \mathrm{d}\mu}
\NewDocumentCommand{\dy}{}{\ \mathrm{d}y}
\NewDocumentCommand{\ind}{}{\chi}
\title[Spectrality of Product Sets with a Perturbed Interval Factor]{Spectrality of Product Sets with a Perturbed Interval Factor}
\author{Aditya Ramabadran, Johannes van Vliet}
\date{Source: arXiv:2508.15159v2 (CC BY 4.0)}
\begin{document}
\maketitle

\noindent\emph{Source and licence.} Spectrality of Product Sets with a Perturbed Interval Factor. Version arXiv:2508.15159v2 (CC BY 4.0), distributed under the Creative Commons Attribution 4.0 International licence, as recorded in the arXiv licence metadata for this exact version and shown on the abstract page https://arxiv.org/abs/2508.15159v2. Full rights notice: licenses/arxiv-2508.15159-CC-BY-4.0.txt.\par\smallskip

\noindent\textit{Editorial scope.} Counted unit: the Proposition whose source label is \texttt{None} in the article (source0.tex lines 199--201, source label \texttt{None}), delivered together with the article's own complete original proof of it (source0.tex lines 203--230, one \texttt{proof} environment of 28 lines). No mathematics is written, repaired or re-derived: statement and proof are the article's own text line for line. Required context retained uncounted: none. Every definition, display and label of those lines is retained unchanged. Omitted from this package and not redistributed: 1--198, 202--202, 231--305. Nothing inside a retained excerpt is omitted or altered: every retained body is delivered exactly as the article's lines are written, so no omission receipt of any kind accompanies this package. Citations: the retained excerpts carry no citation command at all, so no bibliography entry of the article is transcribed and no reference is redistributed. Reference adaptation: none was needed and none was made; every reference inside the delivered unit resolves inside it, so no unresolved reference or citation is produced. Printed slips kept literal: no printed word, symbol, hypothesis or number of the counted statement or of its proof was altered. Layout and class: the article's own class is \texttt{amsart} with packages including hyperref, mathtools, color, geometry, amsmath, amsfonts, amsthm, amssymb, setspace, Tabbing, fancyhdr, lastpage, extramarks, chngpage; the wrapper is the lane's pinned \texttt{amsart} editorial wrapper at 10pt on A4, which restates the theorem-like environments the retained text needs and adds the article's own macro definitions that the retained text needs (\texttt{\textbackslash dmu}, \texttt{\textbackslash dy}, \texttt{\textbackslash essinf}, \texttt{\textbackslash esssup}, \texttt{\textbackslash ind}), with extra wrapper packages (bm, bbm, dsfont, xspace, cancel, mathtools). The delivered numbering is the unit's own. Assets: no retained span contains a figure environment, a graphics-inclusion command, a PDF-page inclusion, a table or a TikZ picture; no image file is copied into this package, the package declares no asset dependency and no image file is redistributed with it. Licence: version arXiv:2508.15159v2 of this article is distributed under the Creative Commons Attribution 4.0 International licence, as recorded in the arXiv licence metadata for this exact version and shown on the abstract page https://arxiv.org/abs/2508.15159v2. A full rights notice is delivered at licenses/arxiv-2508.15159-CC-BY-4.0.txt. Characters: every retained excerpt is pure ASCII, so the delivered engine, XeLaTeX with its default Latin Modern text font, reports no missing character.
\begin{proposition}
Let $m(E)=1$, $\sup E=\esssup E$, $\inf E=\essinf E$, and $E\subset [0,L]$, where $L=3/2-\epsilon$ for some $\epsilon>0$. Suppose $E$ weakly tiles its complement with respect to a locally finite measure $\mu$. Then $\mu=\sum_{k\in\mathbb{Z}\smallsetminus \{0\}}\delta_k$.
\end{proposition}

\begin{proof}
    We will first show that $\mu(-1,1)=0$. Let $K(t) = (\ind_E \ast \ind_{-E})(t) = \int \ind_E(y) \ind_{E}(y-t) \dy = \int \ind_{E \cap (E+t)}(y) \dy = m(E \cap (E+t))$. Then since $\ind_E, \ind_{-E}$ are in $L^2$, $K$ is uniformly continuous. $K$ is also bounded since $m(E \cap (E+t)) \le 1$, and it is compactly supported since $E \subset [0, \frac 32 - \epsilon]$ so for $|t| > \frac 3 2$, $K(t) = 0$. By Lemma 2.1, $K(t) > 0$ for all $t \in [0,1)$. And since Lebesgue measure is translation invariant, $K$ is even and thus $K>0$ on $(-1,1)$. By the Extreme Value Theorem, for each $n$ there exists $m_n>0$ such that $K(t) \ge m_n$ whenever $t \in [-1+\frac{1}{n}, 1 - \frac 1 n]$. 
    
    $E$ weakly tiles its complement with respect to $\mu$, so $\ind_E \ast \mu = \ind_{E^c}$ $m$-almost everywhere.   This means that $\ind_E \ast \mu \ast \ind_{-E} = \ind_{E^c} \ast \ind_{-E}$ a.e. But also, $\ind_E \ast \mu \ast \ind_{-E} = K \ast \mu$ a.e. So $\mu \ast K = \ind_{E^c} \ast \ind_{-E} = (1 - \ind_E) \ast \ind_{-E}$. Then $(1 \ast \ind_{-E})(t) = \int \ind_E(y-t) \dy =m(E)= 1$, so we get $\mu \ast K = 1 - K$ a.e. Since $K$ is uniformly continuous with compact support and $\mu$ is locally finite, the estimate $|(\mu \ast K)(x) - (\mu \ast K)(y)| \le \int |K(x-z) - K(y-z)| \dmu(z)$ tells us that $\mu*K$ is continuous. Since both $1-K$ and $\mu*K$ are continuous and are equal to one another a.e., they are equal everywhere. In particular, we have $(\mu \ast K)(0) = 1 - K(0) = 0$. Then for each $n$, $$0 = \int K(-y) \dmu(y) \ge \int_{[-1 + \frac 1 n, 1-\frac{1}{n}]} K(-y) \dmu(y) \ge m_n \: \mu([-1+\frac 1 n, 1-\frac{1}{n}]),$$ so $\mu([-1 + \frac 1 n, 1-\frac{1}{n}]) = 0$ for all $n$. By continuity from below, we get $\mu((-1, 1)) = 0$. 
    
    Let $r=\sup E=\esssup E$ and let $t=\esssup\left([0,r-1]\cap E^c\right)$. Assume that $t<r-1$. Let $g(x)=\chi_E*\chi_{(1-r,-t)}(x)=m(E\cap((t,r-1)+x))$. Then $g$ is uniformly continuous and $g*\mu$ is continuous by the same reasoning which showed that $K$ and $K*\mu$ were continuous and uniformly continuous, respectively. Also, $g>0$ on $(1,r-t)$ by the minimality of $r$, so for each $n$ there exists $m_n$ such that $g\ge m_n>0$ on $[1+1/n,r-t-1/n]$ by the Extreme Value Theorem. Moreover, because $\chi_E*\mu=1-\chi_E$ almost everywhere, $g*\mu=(1-\chi_E)*\chi_{(1-r,-t)}=r-1-t-g$ almost everywhere. Since $g$ and $g*\mu$ are both continuous, we have equality everywhere, and in particular $g*\mu(0)=r-1-t-g(0)$. And by the maximality of $t$, $g(0)=m(E\cap(t,r-1))=r-t-1$, so $g*\mu(0)=0$. So for any $n$,
    $$0=g*\mu(0)=\int g(-y)d\mu(y)\ge \int_{[t-r+1/n,-1-1/n]}g(-y)d\mu(y)\ge m_n\mu[t-r+1/n,-1-1/n],$$
    which shows that $\mu(t-r,-1)=0$. Furthermore,
    $$\{x\in (t,r-1):-1\in x-E\}=\left((t+1,r)\cap E\right)-1,$$
    and since $m((t+1,r)\cap E)>0$, $\{x\in (t,r-1):-1\in x-E\}$ has positive measure. And since almost all $x\in (t-1,r)$ are contained in $E$ and almost all $x\in E$ satisfy $\mu(x-E)=0$, this shows that $\mu\{-1\}=0$. Thus $\mu(t-r,-1]=\mu(t-r,1)=0$.

    Assume now that $t\leq r-1$. It is still true, then, that $\mu(t-r,1)=0$. By the minimality of $t$, we know that $m(E^c\cap (t-\delta,t))>0$ for any $\delta>0$. This means that we may choose a sequence $y_n$ which increases to $t$ such that $\mu(y_n-E)=1$ for each $n$. Then since $\mu(t-r,1)=0$, we have
    $$\mu(y_n-E)=\mu\left((y_n-E)\cap(-\infty,t-r]\right)\leq \mu(y_n-r,t-r],$$
    and so $\mu\{t-r\}\ge 1$ by continuity from above since $\mu$ is locally finite. 
    
    We know that $t-r\leq -1$, and $m(E\cap[0,1])>0$, so $m(E^c\cap(E+t-r))>0$. Also, $\mu(x-E)=1$ for almost all $x\in E^c$, which means that for some $x\in E^c\cap(E+t-r)$, $\mu(x-E)=1$. Because $x\in E+t-r$, $t-r\in x-E$, which shows that $\mu\{t-r\}\leq\mu(x-E)\leq 1$. Thus $\mu\{t-r\}=1$.

    Let $\mu_1(A)=(\mu+\delta_0-\delta_{t-r})(A+t-r)$, so that $\mu_1$ is a locally finite Borel measure. We have
    $$\chi_E*\mu_1=\tau_{r-t}((\mu+\delta_0)*\chi_E)-\delta_{t-r}(x-E+t-r)=1-\chi_E=\chi_{E^c},$$
    where the third equality holds a.e. due to the fact that $\chi_E*\mu=\chi_{E^c}$ a.e. The reasoning above then shows that $\mu_1(t-r,0)=0$ and that $\mu_1\{t-r\}=1$. Therefore
    $$\mu(2(t-r),t-r)=(\mu+\delta_0-\delta_{t-r})(2(t-r),t-r)=\mu_1(t-r,0)=0,$$
    and
    $$\mu\{2(t-r)\}=(\mu+\delta_0-\delta_{t-r})\{2(t-r)\}=\mu_1\{t-r\}=1.$$
    Continuing in this manner, we find that $\mu|_{(-\infty,0]}=\sum_{j=1}^\infty \delta_{j(t-r)}$. 
    
    Let $E'=r-E$, and let $\mu'(A)=\mu(-A)$. Then $E'$ satisfies the hypotheses of the Proposition and $\mu'$ is a locally finite Borel measure. Also,
    $$\chi_{E'}*\mu'(x)=\mu(-x+r-E)=\chi_E*\mu(-x+r)=\chi_{(E')^c}(x),$$
    with the last equality holding a.e. due to the fact that $\chi_E*\mu=\chi_{E^c}$ a.e. Let $r'=\sup E'=\esssup E'$ and $t'=\esssup[0,r'-1]\cap (E')^c$. Then arguing as we have above, we find that $\mu|_{[0,\infty)}=\sum_{k=1}^\infty \delta_{k(r'-t')}$. This means that $\mu+\delta_0=\sum_{\lambda\in \Lambda}\delta_\lambda$, where $\Lambda$ is a discrete set and $|\lambda-\lambda'|\ge 1$ whenever $\lambda\neq\lambda'$. Therefore $[0,1]+\Lambda$ is a packing. And since $E+\Lambda$ is a tiling, $\Lambda$ must be a tiling set for $[0,1]$ by Theorem 2.2. Because $0\in\Lambda$, we must have $\Lambda=\mathbb{Z}$. It follows that $\mu=\sum_{k\in\mathbb{Z}\smallsetminus\{0\}}\delta_k$. \end{proof}

\end{document}
